% !PRINTED 552
\weberpage{552}
% running head: 552 Siebzehnter Abschnitt. \S.\,152.
%<<CONT
divisible number can be represented as a product of ideal
numbers. All these questions are answered through the mediation
of the theory of functionals (or the theory of ideals), if one
assigns to the ideal number $(\alpha, \beta)$ the functional $\varphi$, and they can
scarcely be decided in any simpler way. The introduction
of the system of ideal numbers then offers no further
essential advantages, and one can just as well regard the
functionals themselves, or the ideals, as elements of the
manifold that is to be newly introduced.

\begin{center}
\S.\,152.

\gesperrt{Equivalence.}
\end{center}

We now call two functionals that differ only by a
unit factor, even when they are fractional,
\gesperrt{associated}, and set up the following definition:

1. \gesperrt{Two integral or fractional functionals $\varphi, \psi$
in the field $\Omega$ are called equivalent if their
quotient $\varphi : \psi$ is associated with a number.}

Thus the two functionals $\varphi$ and $\psi$ are called equivalent
if there exist a unit $\varepsilon$ and a number $\alpha$ in $\Omega$ such that

\weq{1}{\frac{\varphi}{\psi} = \alpha\varepsilon}

holds. If $\varphi$ is equivalent with $\psi$ and with $\psi_1$, it follows from (1) that

\[
\frac{\varphi}{\psi} = \alpha\varepsilon, \qquad \frac{\varphi}{\psi_1} = \alpha_1\varepsilon_1,
\]

consequently

\[
\frac{\psi}{\psi_1} = \frac{\alpha_1}{\alpha} \frac{\varepsilon_1}{\varepsilon},
\]

and since $\alpha_1 : \alpha$ is a number and $\varepsilon_1 : \varepsilon$ a unit, the
first theorem follows:

2. \gesperrt{Two functionals that are equivalent to a third
 are also equivalent to one another
.}

If, in accordance with this, one divides all the functionals of
the field $\Omega$ into \gesperrt{classes}, by casting two functionals into the
same or into different classes according as they are equivalent
or not, it results that two of these classes which have a single
%CONT>>
